\documentclass{article}
\usepackage[T1]{fontenc}
\usepackage{lmodern}
\usepackage{graphicx}
\usepackage{amsmath,amssymb}
\usepackage[a4paper,margin=2cm]{geometry}
\usepackage[absolute,overlay]{textpos}
\setlength{\TPHorizModule}{1mm}
\setlength{\TPVertModule}{1mm}
\usepackage[indonesian]{babel}
\usepackage{hyperref}
\setlength{\emergencystretch}{2em}
% unit: Halaman 55

\begin{document}

% === Halaman 55 (python/native) ===

Kriptografi di Eropa dari Zaman Renaisans sampai Abad 19

• Zaman renaisans → abad pertengahan (abad 15-16)

• Cipher terkenal pada abad pertengahan hingga abad 19: 

   1. Vigenere Cipher 

       Dipublikasikan oleh diplomat Perancis bernama Blaise de Vigenere 

pada tahun 1586.

    2. Playfair Cipher

        Dipromosikan oleh diplomat Inggris, Lord Playfair, meskipun  

penemu aslinya adalah Charles Wheastone pada tahun 1854.

55

\end{document}
